\section{Tooling 与自动化}

\subsection{自动化流水线}

Offline RL 需要把 data -> training -> eval -> deploy 串成流水线。典型组件：

\begin{itemize}
    \item Data versioning：每次训练都指向数据、reward log、behavior policy ID。
    \item Training orchestration：Pipeline.com、Kubeflow pipeline 触发 training job。
    \item Evaluation orchestration：自动化地从 replay buffer 取 key states 运行 benchmark。
    \item Deployment gate：若 metrics 异常，自动回滚并通知 owner。
\end{itemize}

\begin{warningbox}{流水线中断会导致 drift}
一旦 pipeline 的任一环节（data/eval/deploy）缺失，后续变化就无法追踪。必须强制所有 code change 通过 pipeline 生成的 artifact，避免“手动跑”。
\end{warningbox}

\subsection{支持工具与集成}

讲者列举了几款与 offline RL 结合良好的工具：

\begin{itemize}
    \item \textbf{Data-versioned storage}（Feathr, Deep Lake）存储 dataset metadata。
    \item \textbf{Experiment tracking}（Weights & Biases）记录 Q variance、KL drift。
    \item \textbf{Policy registry}：记录 policy hash, deploy timestamp, sign-off, governance doc link。
    \item \textbf{Replay simulator}：用 deterministic seeds replay 训练轨迹并记录 divergences。
\end{itemize}

\subsection{本章小结}

Automation pipeline 不只是流水线，而是把 offline RL 的各个 artifact 变成 traceable product。实时 logging + tooling 集成让 governance 更容易。

